\begin{python0}
from solutions import *; clear()
\end{python0}
\sectionthree{Lazy allocation}

There are times when you might want to delay allocating memory (or any
resource) because it might not be needed till later.

You can assign any pointer the value of NULL.

NULL (or the NULL pointer) is usually integer 0 (or something similar
... like \texttt{(void *)0}). NULL is frequently used as a value to
indicate that a pointer has not been allocated memory.

\begin{consolethree}[escapeinside=||]
#include <cstdlib> // or #include <cstddef>
#include <iostream>

// NULL can be the integer 0 or it can be
// ((void *) 0)
class C{};

int main()
{
    // NULL can be assigned to pointers of any type:
    int * p = NULL;
    char * q = NULL;
    double * r = NULL;
    bool * s = NULL;
    C * t = NULL;
    std::cout << p << '\n';
    std::cout << (int)NULL << '\n';
    return 0;
}
\end{consolethree}

In this case you can do the following: In the constructor, do NOT
allocate memory but set the \EMPHASIZE{relevant pointer} to NULL.

Example:

\begin{consolethree}[escapeinside=||]
...
class LazyIntPointer
{
public:
    LazyIntPointer()
        : p_(NULL)
    {}
    ...
private:
    int * p_;
};
...
\end{consolethree}

When do you allocate memory? Only when you need to use the value that
p\_ points to:

\begin{consolethree}[escapeinside=||]
...
class LazyIntPointer
{
public:
    ...
    int & operator*()
    {
        // Allocate memory only when p_ is NULL
        if (p_ == NULL)
        {
            p_ = new int;
        }
        return *p_;
    }
    ...
};
...
\end{consolethree}

And in the destructor, only deallocate when \texttt{p\_} is not NULL,
i.e., only when \texttt{p\_} has actually been allocated memory:

\begin{consolethree}[escapeinside=||]
...
class LazyIntPointer
{
public:
    ...
    ~LazyIntPointer()
    {
        // De-allocate memory only when p is not NULL
        if (p_ != NULL)
        {
            delete p_;
        }
    }
    ...
};
...
\end{consolethree}

This is a very common technique to save on memory usage. (There are many
other techniques.)


